\section{Related work}
% === ROUND 2 --- §2 goals & contents ======================================
% Section goal: pin the empty cell PRECISELY and set the measurement-validity bar
% this work is judged against. Treatment: FULL REWRITE to the §1 standard
% (claims/citations/table kept; prose rewritten under the doctrine + readability lessons).
%  2.1 grown/measured -- GOAL: the signal exists & is non-trivial but never
%      intervenes. CONTAINS: scale/layer/saturation findings; takeaway = silent on usability.
%  2.2 grown/optimised -- GOAL: the bare inversion is not novel, but no one compresses.
%      CONTAINS: speech+text brain-tuning, the causal result; the gap (all GROW).
%  2.3 shrunk/measured -- GOAL: the closest prior work, scoped exactly. CONTAINS:
%      oota2026 (narrow quantization survival; pruning/KD untested); why KD is lossier; DPI framing.
%  2.4 empty cell + weak precedents -- GOAL: the cell is unoccupied; adjacent
%      precedents weaken. CONTAINS: the empty cell; vision/audio precedents + the missing matched control.
%  2.5 measurement bar + surrogate -- GOAL: the standard + why a surrogate is future
%      work. CONTAINS: leakage/artifact cautions; the surrogate direction is unsettled.
% ===========================================================================

This thesis sits in the one empty cell of a two-by-two grid: grow or shrink the model, and measure or optimize alignment (Table~\ref{tab:cells}).
We review the three occupied cells in turn, then the measurement standard this work adopts.

\begin{table}[H]
  \centering
  \footnotesize
  \caption{The field on two axes: what an intervention does to the parameter budget, and how it treats brain alignment. Three cells are occupied; the fourth, shrinking a model while alignment is in the objective, is empty in the 2026 literature and is where this work sits.}
  \label{tab:cells}
  \resizebox{\textwidth}{!}{%
  \begin{tabular}{lll}
    \toprule
                       & Alignment measured & Alignment optimized \\
    \midrule
    Model grown        & classical brain-score \autocite{oota2023,gao2024,alkhamissi2025} & brain-tuning \autocite{moussa2025,bilgin2026,merlin2026} \\
    Model shrunk       & compression robustness \autocite{oota2026} & \textbf{this thesis (empty in the literature)} \\
    \bottomrule
  \end{tabular}}
\end{table}

\subsection{The classical brain-score literature}

The largest body of work fits encoding models to frozen, pretrained networks, then reads off how alignment varies with architecture, scale, and training.
Three findings recur.
Alignment lives in the structured middle layers, where syntax drives the layerwise trend and semantics is more region-specific \autocite{oota2023}.
It grows with model scale and saturates around a few billion parameters, and at a matched parameter count scale moves it more than instruction-tuning does, which is why a base model is the right teacher to align against \autocite{gao2024}.
Beyond that scale the trend can flatten or reverse, as the largest models outgrow the human language network and track formal competence more than behavior \autocite{alkhamissi2025}.
This line of work shows the signal exists and is structured, but it never intervenes through alignment, so it cannot say whether alignment is usable.

\subsection{Brain-tuning}

A second line puts alignment in the objective of a fine-tuned or grown model, across modalities.
For speech, a voxelwise brain loss fine-tunes compact encoders and improves both alignment and several downstream tasks, with a multi-participant variant reporting further gains at higher data efficiency \autocite{moussa2025,moussa2025b}.
For text, training GPT-2 and LLaMA-2 with a cosine brain loss beats text-only baselines, though the gains are mixed and the splits are shuffled with no brain-shuffled control, so they would not clear the measurement bar this work adopts \autocite{bilgin2026}.
The strongest evidence is causal.
Training a text model to be \emph{less} brain-aligned, at fixed perplexity, makes it worse across a suite of more than two hundred tasks, while training it to be \emph{more} aligned improves performance on the same suite \autocite{merlin2026}.
Alignment is therefore doing causal work, not just riding along with a good model.
The bare inverted idea is not novel, but no study here compresses a model: every method grows or holds the parameter budget, leaving the compression question open.

\subsection{Alignment under compression}

The closest prior work to this thesis is the compression study of \textcite{oota2026}, which tracks how alignment changes as a model is compressed.
They find that alignment survives only certain post-training quantizers on models of at least a few billion parameters; other quantizers, smaller models, and pruning all degrade it, and distillation is never tested.
That study measures alignment under compression but never optimizes it. It reports how quantization and pruning move the score, and never makes alignment a training target.
Distillation is a different and lossier operation.
It refits a smaller model from scratch against the teacher, the step most likely to destroy the middle-layer structure where alignment lives.
So alignment surviving quantization, a gentle reshaping of fixed weights, does not imply it survives distillation.
Whether distillation preserves alignment is therefore an open empirical question: how much does it lose by default, and can a brain objective recover the rest?
We develop the information-theoretic bound behind this in Section~\ref{sec:methods}.

\subsection{The empty cell}

No prior work occupies the empty cell.
Two nearby precedents, one in vision and one in audio, are sometimes cited as indirect support for a brain regularizer when data is scarce, but neither isolates a brain-specific effect.
\textcite{pirlot2022} regularize a vision network with neural data and improve both accuracy and adversarial robustness, yet a shuffled-label control reproduces the accuracy gain, so only the robustness benefit depends on real neural structure.
\textcite{freteault2025} fine-tune a small audio network on fMRI and improve most downstream tasks, but without an interaction test or a matched non-brain control, the low-data advantage is not isolated.
Neither study is in the distillation regime this work targets, so both inform our design without testing the claim.

\subsection{The measurement bar}

A separate line of work warns that alignment scores are easy to inflate and over-read, and it sets the standard this work adopts.
Shuffled splits leak through the temporal autocorrelation of fMRI, low-level features like sentence length carry much of the apparent signal, and some published scores fail basic predictive and reliability checks \autocite{hadidi2026,jia2026}.
The same headline score can reflect real semantics in one model and shallow phonology in another, so the number alone does not pin down the mechanism \autocite{oota2024}.
A differentiable, fMRI-free surrogate would go further, but the literature does not agree on its direction: local intrinsic dimension correlates with alignment in opposite ways across language and vision, and similarity-based surrogates track width and depth as much as alignment \autocite{cheng2026,yu2026,groger2026}.
We therefore leave a surrogate to future work, since real fMRI must first establish which way to push it.
